% =============================================================================
\section{The Transformer Architecture}
\label{sec:transformer}
% =============================================================================

\subsection{Motivation and Core Idea}

Both LSTM variants process text token by token. The root cause of their
limitations is the sequential recurrence: dependency on the previous hidden
state forces one-at-a-time computation and creates an information bottleneck for
long-range dependencies.

The \textbf{Transformer}~\cite{vaswani2017attention} (Vaswani et al., 2017)
discards recurrence entirely. Its central mechanism is \textbf{self-attention}:
every token directly attends to every other token in the sequence
\emph{simultaneously}, computing a weighted combination of all position
representations in a single operation.

\subsection{Self-Attention Mechanism}
\label{subsec:selfattn}

Given an input sequence of $n$ tokens represented as a matrix
$X \in \R^{n \times d_{\text{model}}}$, three linear projections produce query,
key, and value matrices:
\begin{equation}
  Q = X W^Q, \quad K = X W^K, \quad V = X W^V
  \label{eq:qkv}
\end{equation}
where $W^Q, W^K, W^V \in \R^{d_{\text{model}} \times d_k}$.
The \textbf{scaled dot-product attention} output is:
\begin{equation}
  \text{Attention}(Q, K, V)
    = \softmax\!\left(\frac{QK^\top}{\sqrt{d_k}}\right) V
  \label{eq:attention}
\end{equation}

The scaling factor $1/\sqrt{d_k}$ prevents the dot products from growing large
in magnitude, which would push the softmax into regions of extremely small
gradients. Each output row is a weighted sum of all value vectors, where the
weights reflect how relevant each token is to the current query position.

\subsection{Multi-Head Attention}

Rather than performing a single attention operation, the Transformer applies $h$
\emph{parallel} attention heads, each learning a different projection:
\begin{equation}
  \text{MultiHead}(Q,K,V)
    = \text{Concat}(\text{head}_1, \ldots, \text{head}_h)\, W^O
  \label{eq:multihead}
\end{equation}
where $\text{head}_i = \text{Attention}(QW_i^Q,\, KW_i^K,\, VW_i^V)$ and
$W^O \in \R^{hd_k \times d_{\text{model}}}$. Multiple heads allow the model
to capture different types of relationships simultaneously --- \textit{e.g.}\
syntactic agreement, coreference, and semantic similarity --- from the same
input.

% ── Self-attention illustration ───────────────────────────────────────────────
\begin{figure}[H]
\centering
\begin{tikzpicture}[>=Stealth, scale=1.02, transform shape]
  \def\words{{"The","animal","didn't","cross","the","street","because","it"}}
  \foreach \i in {0,...,7}{
    \pgfmathsetmacro{\lbl}{\words[\i]}
    \node[draw=NavyBlue!85, fill=NavyBlue!20, rounded corners=3pt,
          minimum width=1.25cm, minimum height=0.68cm,
          font=\footnotesize\bfseries, align=center, text=black]
          (tok\i) at ({-5.4 + \i*1.55}, 0) {\lbl};
  }
  \foreach \i in {0,1,2,3,4,5,6}{
    \pgfmathsetmacro{\opac}{0.18 + \i*0.08}
    \draw[cyan!65!blue, opacity=\opac, line width={\opac*6pt}]
      (tok7.south) .. controls +(0,-0.9) and +(0,-0.9) .. (tok\i.south);
  }
  \draw[red!80!black, very thick, line width=2.8pt]
    (tok7.south) .. controls +(0,-1.35) and +(0,-1.35) .. (tok1.south);
  \node[font=\normalsize\bfseries, text=red!80!black] at (-1.0,-1.65)
    {highest attention: ``it'' $\rightarrow$ ``animal''};
  \node[font=\small, text=black!85, align=center] at (0,-2.45)
    {Each token attends to \emph{all} others simultaneously (in parallel).
     Long-range links resolved in one step.};
\end{tikzpicture}
\caption{Self-attention in action. The pronoun \textit{``it''} attends most
  strongly to \textit{``animal''}, resolving the coreference in a single layer
  without traversing a chain of timesteps. Thicker arcs indicate higher
  attention weights.}
\label{fig:selfattn}
\end{figure}

\subsection{Positional Encoding}

Since self-attention is permutation-invariant (Eq.~\ref{eq:attention} has no
notion of token order), the Transformer injects \textbf{positional
encodings}~\cite{vaswani2017attention} into the input embeddings:
\begin{equation}
  \text{PE}_{(pos,\,2i)}   = \sin\!\left(\frac{pos}{10000^{2i/d}}\right),
  \qquad
  \text{PE}_{(pos,\,2i+1)} = \cos\!\left(\frac{pos}{10000^{2i/d}}\right)
  \label{eq:posenc}
\end{equation}
These fixed sinusoidal functions encode absolute position and allow the model to
represent relative distances between tokens.

\subsection{Encoder and Decoder Stacks}

The original Transformer~\cite{vaswani2017attention} is an encoder--decoder
model. Each encoder layer consists of (i)~multi-head self-attention and
(ii)~a position-wise feed-forward network (FFN), each sublayer wrapped with a
residual connection and Layer Normalisation. The decoder additionally contains
a \emph{cross-attention} sublayer that queries the encoder output.

\begin{figure}[H]
  \centering
  \includegraphics[width=0.75\textwidth]{og_transformer.png}
  \caption{The original Transformer architecture~\cite{vaswani2017attention}.
    The encoder (left) produces contextual embeddings from the source sequence.
    The decoder (right) generates the target sequence auto-regressively, using
    masked self-attention and cross-attention over the encoder output.}
  \label{fig:transformer}
\end{figure}

\subsection{Why the Transformer Surpasses the LSTM}

Table~\ref{tab:lstm_vs_transformer} summarises the key differences between the
two architectures.

\begin{table}[H]
\centering
\caption{LSTM vs.\ Transformer: a direct comparison.}
\label{tab:lstm_vs_transformer}
\begin{tabularx}{0.92\textwidth}{lXX}
  \toprule
  \textbf{Property} & \textbf{LSTM / BiLSTM} & \textbf{Transformer} \\
  \midrule
  Parallelism
    & Sequential; one step per token
    & Fully parallel over all tokens \\
  Long-range dependencies
    & Degrade with distance
    & Constant path length between any two positions \\
  Bidirectionality
    & Separate forward/backward passes
    & Joint left+right context within every attention layer \\
  Scalability
    & Hard to scale to very long sequences
    & Scales with hardware via parallelism \\
  \bottomrule
\end{tabularx}
\end{table}
